% vim: tw=78 ai spell spelllang=cs

\documentclass{article}

\usepackage{graphicx}
\usepackage{hyperref}
\usepackage[utf8x]{inputenc}
\usepackage[czech]{babel}
\usepackage[T1]{fontenc}
\usepackage{xspace}

\pagestyle{empty}

\newcommand{\topinfo}[4]{
\begin{flushright}
\today
\end{flushright}
\vspace{1em}
\begin{center}
\Large\textbf{Posudek #1 na #2 práci \\
\normalsize%
\textit{#3:} #4}
\end{center}
}
\newcommand{\botinfo}{
\vspace{3em}
\begin{flushright}
Jiří Slabý
\end{flushright}
}


\newcommand{\symb}{\textsc{Symbiotic}\xspace}

\begin{document}
\topinfo{vedoucího}{bakalářskou}{Marek Chalupa}{Analýzy ukazatelů pro
Symbiotic}

Dle zadání měla bakalářská práce nastudovat analýzy ukazatelů, implementovat
Shapiro-Horwitzové a porovnat s existující Andersenovou v rámci nástroje
\symb. Text práce je na 35 stranách strukturovaný do 7 kapitol.  V úvodu autor
krátce popisuje motivaci pro práci.  Následující kapitola popisuje 3 analýzy
ukazatelů kódu včetně základních pojmů, teorie a jejich porovnání.
Následuje kapitola s popisem nástroje \symb a jeho analýz, což je základem pro
další kapitoly: implementaci analýzy Shapiro-Horwitzové a úvahu o rozšíření na
pole. Autor pokračuje v 6.\ kapitole měřením a porovnáním implementací. Práci
ukončuje závěr.  \textbf{Rozsah} uvedených kapitol podle mého názoru převyšuje
standardní nároky na bakalářské práce i vzhledem k zadání.

\textbf{Obsahově} dle mého názoru dostatečně popisuje teoretické základy
práce, návrh i implementaci. Implementace a její rychlost jsou ověřeny na
několika reálných programech.

Po \textbf{jazykové} stránce je práce na dobré úrovni. Je zde jen několik
překlepů (např.\ ,,field sensitiviy'' na str.\ 4) a nekonzistencí (např.\
,,ořezat'' vs.\ ,,prořezat'' kód).

\textbf{Praktická část} práce je v přiloženém \texttt{tar.gz} souboru a také na
serveru \url{http://github.com}. Jelikož se jedná o rozšíření již existujícího
nástroje \symb, je v práci uvedeno, co přesně bylo změněno a přidáno.  Tento
nově přidaný kód je čitelný a jeví se být plně funkční.

Na závěr konstatuji, že práce \textbf{splňuje zadání a požadavky na
bakalářskou práci}. Navrhuji práci \textbf{uznat} jako bakalářskou a i přes
výše uvedené okolnosti doporučuji hodnocení stupněm A.

\botinfo

\end{document}
